\documentclass[11pt,a4paper]{coverletter}

\senderblock
  {Radheem Bin Razi}
  {Distributed Systems \& Cloud-Native Engineer}
  {Ilmenau, Germany \textbar\ \href{mailto:sheikh.radheem@gmail.com}{sheikh.radheem@gmail.com} \textbar\ \href{https://radheem.github.io/my-cv}{portfolio}}

\begin{document}

%% ===== ENGLISH =====
\recipient{Remerge}{Hiring Team}{}

\opening{Dear Hiring Team,}

Your Performance team's view of ML is exactly how I work: feature pipelines, online serving, and monitoring are part of the job, not someone else's problem. I have built that full path with Python and TensorFlow, and I am comfortable reading Go to navigate online-serving services — which is what drew me to this Machine Learning Engineer role at Remerge.

Most directly, my O-RAN AIML research project is an end-to-end ML platform: I designed a config-driven Python SDK that runs Kubeflow training and retraining pipelines and serves models online with KServe, backed by a feature store and TensorFlow/Keras + scikit-learn prediction models for per-cell throughput. It took raw telemetry through feature definitions, training, evaluation, and serving — the same arc as taking a bidding model from raw data to online inference. The work was graded 1.0 (A). Earlier, at Seed Labs, I owned data problems end to end, cleaning and analyzing data and building scikit-learn and TensorFlow models that delivered concrete insight to stakeholders.

What I would bring to the bidding engine is the production-systems side that makes ML hold up under load. At Bluefin Exchange I was the design authority for core services of a high-throughput crypto exchange — a real-time, latency-sensitive system on Go, Kafka, PostgreSQL, and Kubernetes — where I reviewed designs and code and kept releases reliable. That is the "honest about cost and scale" mindset your hybrid train-in-cloud, serve-on-own-hardware setup demands. I will be candid that I have not built RTB systems specifically, but I have built the high-throughput serving and streaming foundations they run on, and I lean on a pragmatic, data-driven read of what a problem actually needs. I also work daily with AI coding tools to keep experiment velocity high.

I am based in Ilmenau, Germany, available full-time now (open to working-student/part-time), and Blue Card-ready — so you only issue the contract, with no sponsorship overhead. Remerge's remote-first setup fits well, and I can relocate within Germany in about two to three weeks if useful. I would welcome the chance to talk.

\closing{Sincerely,}{Radheem Bin Razi}

\clearpage
\selectlanguage{ngerman}

%% ===== DEUTSCH =====
\recipient{Remerge}{Personalabteilung}{}

\opening{Sehr geehrtes Hiring-Team,}

Das ML-Verständnis Ihres Performance-Teams entspricht genau meiner Arbeitsweise: Feature-Pipelines, Online-Serving und Monitoring sind Teil der Aufgabe und nicht das Problem von jemand anderem. Ich habe diesen gesamten Weg mit Python und TensorFlow aufgebaut und kann Go lesen, um mich durch Online-Serving-Services zu bewegen — und genau das hat mich an dieser Position als Machine Learning Engineer bei Remerge gereizt.

Am unmittelbarsten ist mein O-RAN-AIML-Forschungsprojekt eine End-to-End-ML-Plattform: Ich habe ein konfigurationsgetriebenes Python-SDK entworfen, das Kubeflow-Trainings- und Retraining-Pipelines ausführt und Modelle online mit KServe bereitstellt, gestützt auf einen Feature-Store und TensorFlow/Keras- + scikit-learn-Vorhersagemodelle für den Durchsatz pro Zelle. Es führte Rohtelemetrie durch Feature-Definitionen, Training, Evaluierung und Serving — derselbe Bogen wie der Weg eines Bidding-Modells von Rohdaten bis zur Online-Inferenz. Die Arbeit wurde mit 1,0 (A) bewertet. Zuvor habe ich bei Seed Labs Datenprobleme von Anfang bis Ende verantwortet, Daten bereinigt und analysiert und scikit-learn- und TensorFlow-Modelle erstellt, die Stakeholdern konkrete Erkenntnisse lieferten.

Was ich in die Bidding-Engine einbringen würde, ist die Produktionssystem-Seite, die ML unter Last bestehen lässt. Bei Bluefin Exchange war ich die Design-Autorität für die Kernservices einer Hochdurchsatz-Kryptobörse — ein echtzeitfähiges, latenzkritisches System auf Go, Kafka, PostgreSQL und Kubernetes — wo ich Designs und Code geprüft und Releases zuverlässig gehalten habe. Das ist die Denkweise des "ehrlichen Umgangs mit Kosten und Skalierung", die Ihr hybrides Setup aus Training in der Cloud und Serving auf eigener Hardware erfordert. Ich werde offen sagen, dass ich nicht speziell RTB-Systeme gebaut habe, aber ich habe die Hochdurchsatz-Serving- und Streaming-Grundlagen aufgebaut, auf denen sie laufen, und ich verlasse mich auf eine pragmatische, datengetriebene Einschätzung dessen, was ein Problem tatsächlich erfordert. Außerdem arbeite ich täglich mit KI-Coding-Tools, um die Experimentiergeschwindigkeit hoch zu halten.

Ich bin in Ilmenau, Deutschland, ansässig, ab sofort in Vollzeit verfügbar (offen für Werkstudenten-/Teilzeit) und Blue-Card-bereit — Sie stellen also nur den Vertrag aus, ohne Sponsoring-Aufwand. Das Remote-First-Setup von Remerge passt gut, und ich kann bei Bedarf innerhalb von etwa zwei bis drei Wochen innerhalb Deutschlands umziehen. Über die Gelegenheit zu einem Gespräch würde ich mich freuen.

\closing{Mit freundlichen Grüßen,}{Radheem Bin Razi}

\end{document}
